\section{System Overview \& LLM Selection}
\subsection{Problem Statement}
The School of Computer Science manages a wide range of shared physical spaces, including classrooms, computer laboratories, project laboratories, meeting rooms, auditoriums, and student workspaces. These spaces are used for teaching activities, seminars, examinations, workshops, research, student projects, and academic events. Because many users need access to the same spaces, the School must ensure that room usage is organized, fair, and efficient.

At present, the booking process is handled manually. Lecturers, teaching assistants, students, and staff usually contact the school office or facility staff by email, phone, or in person. Facility staff then check spreadsheets or shared calendars to verify room availability, user eligibility, required facilities, and maintenance status. This manual process becomes increasingly difficult to manage as the number of booking requests grows. It is time-consuming, error-prone, and may lead to conflicting bookings, incomplete records, and inefficient space utilization.

The School therefore needs a database system to support the management of shared campus spaces in a more structured and reliable way. The system should store information about users, bookable spaces, available facilities, booking requests, approval decisions, check-in and completion records, maintenance activities, and historical usage. It must also enforce important business rules, such as preventing overlapping approved bookings for the same space, blocking bookings for spaces that are under maintenance or closed, and recording accountability information for approvals and maintenance actions.

The main goal of the proposed system is to help the School manage shared campus spaces fairly and efficiently, reduce booking conflicts, prevent the use of unavailable spaces, and preserve complete usage history for future reference and decision-making.

\subsection{LLM Models Utilized}
\begin{itemize}
    \item \textbf{Big Pickle:} Evaluated via the OpenCode Zen platform, this model is specifically optimized for automated coding agents. Featuring an extensive context window of approximately 200k tokens, it serves as an ideal tool for rapid initial prototyping. However, its primary limitation lies in the absence of advanced reasoning features.
    
    \item \textbf{DeepSeek V4 Flash Free:} Accessed through the OpenCode platform, this model was deployed to facilitate technical drafting, content refinement, and structural editing (e.g., database design documentation). It demonstrates high efficiency in generating structured explanations and concise academic prose, though it necessitates human oversight to ensure rigorous factual accuracy and alignment with project specifications.
\end{itemize}